\section{Refutable, not confirmable}
\label{sec:argument}

A scalp recording at $D$ electrodes is a projection of far more cortical and non-cortical current sources, smeared by volume conduction \citep{nunez2006electric,michel2019eeg}.
For any observed $x$, infinitely many source configurations are compatible with it, and source-localisation methods restore uniqueness only relative to a chosen prior.
Testing a recovery therefore needs a reference that is independent of the method.
The candidates on an observational recording are inverse solutions, behavioural decoding and concurrent recordings in another modality, and each encodes a model of its own: agreement between two model-based unmixings shows that the two models agree, a decoder absorbs any error in the components that is consistent with the readout, and intracranial contacts index cortical patches rather than the latent sources a nonlinear-ICA theorem speaks about.
A component set that scores well on any of these can be a faithful recovery, a confound or a self-consistent artefact.

The theorems nevertheless leave something testable.
Time-contrastive learning (TCL; \citealp{hyvarinen2016tcl}) requires that the segment-wise modulation of the sources spans the source space, i.e.\ that the modulation matrix has full column rank (Appendix~\ref{app:tcl}); iVAE \citep{khemakhem2020ivae} requires the analogous invertibility for its auxiliary variable; for second-order linear separation the exact condition is that no two sources have proportional variance profiles \citep{pham2001blind,kleinsteuber2013uniqueness}; permutation-contrastive learning (PCL; \citealp{hyvarinen2017pcl}) requires non-Gaussian innovations.
If all sources always get louder together, the modulation matrix has rank one and only the total power is identified (Appendix~\ref{app:tcl}).
This is visible in the segment covariances of the recording itself, with no ground truth, and it refutes the claim.
What no recording can supply is the converse.
A segment classifier that predicts segment identity from slow drift, or a contrastive objective fed by the autocorrelation of $1/f$ noise, satisfies the method's objective as well as the sources would, and the recording offers no quantity that separates the two.
This is Popper's asymmetry in its original form, data can refute a hypothesis but cannot verify it \citep{popper1959logic}, applied to a claim that is routinely reported as verified.
The rest of the paper measures both halves.
